% IFN680 Project 7 technical report.
% Generated by tools/build_report.py from tools/report_template.tex. Every number below is
% printed by main_report.ipynb, except the design settings and the prior-work findings that
% tools/verify.py lists.
\documentclass[10pt,a4paper,twocolumn]{article}

\usepackage[margin=1.5cm,top=0.9cm,bottom=0.8cm]{geometry}
\usepackage{lmodern}
\usepackage[T1]{fontenc}
\usepackage[utf8]{inputenc}
\usepackage{microtype}
\usepackage{graphicx}
\usepackage{booktabs}
\usepackage{amsmath}
\usepackage[font=small,labelfont=bf,skip=3pt]{caption}
\usepackage{titlesec}
% Loaded last, as hyperref requires; it only fills the PDF's title and author fields.
\usepackage[hidelinks,pdftitle={Project 7: Improving DDPM},
            pdfauthor={Karan Rooprai, Nhu Hieu Nguyen}]{hyperref}

\titlespacing*{\section}{0pt}{4pt}{0.5pt}
\titleformat{\section}{\normalfont\bfseries}{\thesection.}{0.5em}{}
\setlength{\parskip}{0pt}
\setlength{\parindent}{1em}
\setlength{\columnsep}{0.6cm}
\setlength{\textfloatsep}{4pt}
\setlength{\dbltextfloatsep}{5pt}
\setlength{\intextsep}{3pt}
\setlength{\floatsep}{3pt plus 1pt minus 1pt}
\setlength{\tabcolsep}{4pt}
\pagestyle{empty}
\linespread{0.97}
% Never split a hyphenated word across a column or page break.
\brokenpenalty=10000

\renewcommand{\topfraction}{0.94}
\renewcommand{\dbltopfraction}{0.94}
\renewcommand{\bottomfraction}{0.6}
\renewcommand{\textfraction}{0.06}
\renewcommand{\floatpagefraction}{0.75}

\begin{document}

% The banner sits inside the \twocolumn optional argument, which is the only way a page-wide
% figure can appear on page 1 in the article class.
\twocolumn[{%
\begin{center}
{\large\bfseries Project 7: Improving DDPM. Does a Cosine Schedule Let a Diffusion Model Sample in
Fewer Steps?}\\[2pt]
{\bfseries Group 4}
\end{center}
\vspace{-2pt}
Student 1: Karan Rooprai \hfill Student 1 ID: n12498122 \\
Student 2: Nhu Hieu Nguyen \hfill Student 2 ID: n12194778\par
\vspace{4pt}
% \centering rather than a center environment, whose \topsep would add space above the figure.
{\centering\includegraphics[width=0.80\textwidth]{figures/figure_1_grids.pdf}\par}
\vspace{2pt}
{\small\textbf{Figure 1:} Twelve samples per digit from the linear-schedule (left) and
cosine-schedule (right) models at seed 0 and 1000 steps; images in the same place start from the
\emph{same} noise. The classifier reads 0.9583 (left) and
0.9250 (right) as the intended digit.\par}
\vspace{6pt}
}]

% A page-wide table can only float to the top of a later page, so it is declared here, on page 1,
% to appear at the top of page 2.
\begin{table*}[t]
\centering
\footnotesize
\begin{tabular}{@{}r rrll rrll@{}}
\toprule
& \multicolumn{4}{c}{DDPM sampler} & \multicolumn{4}{c}{DDIM sampler} \\
\cmidrule(lr){2-5}\cmidrule(l){6-9}
$K$ & Linear & Cosine & $\Delta$ [95\% CI] & $\Delta$, seeds 1, 2 & Linear & Cosine &
$\Delta$ [95\% CI] & $\Delta$, seeds 1, 2 \\
\midrule
10 & 186.26 & \textbf{85.35} & $-$100.91 [$-$111.11, $-$89.05] & $-$125.94, $-$120.97 & 49.58 & 106.91 & $+$57.33 [$+$44.97, $+$69.68] & $+$3.86, $-$4.93 \\
20 & 61.77 & \textbf{26.55} & $-$35.22 [$-$41.54, $-$28.39] & $-$13.18, $-$29.46 & 33.89 & 40.78 & $+$6.89 [$-$0.42, $+$14.88] & $+$16.94, $-$4.85 \\
50 & 26.55 & 17.25 & $-$9.30 [$-$13.84, $-$4.54] & $+$14.59, $-$2.90 & 28.48 & 27.31 & $-$1.17 [$-$7.22, $+$5.27] & $+$25.26, $-$2.12 \\
100 & 18.67 & 15.61 & $-$3.06 [$-$7.57, $+$1.53] & $+$17.78, $+$2.57 & 27.28 & 25.13 & $-$2.14 [$-$7.76, $+$3.97] & $+$26.75, $-$1.24 \\
250 & 15.27 & 15.58 & $+$0.31 [$-$3.64, $+$4.64] & $+$22.65, $+$1.28 & 26.80 & 24.09 & $-$2.71 [$-$8.06, $+$3.21] & $+$27.35, $-$0.82 \\
1000 & 17.12 & 13.50 & $-$3.61 [$-$7.70, $+$0.74] & n/a & n/a & n/a & n/a & n/a \\
\bottomrule
\end{tabular}
\caption{Fr\'echet distance to real test digits (lower is better; real-data floor 7.65) at
$K$ sampling steps, seed-0 models, 1,000 samples each. $\Delta$ is cosine minus linear with a
paired bootstrap interval, then at seeds 1, 2. Bold: better by the fixed rule.}
\end{table*}

\section{Introduction}
A denoising diffusion probabilistic model (DDPM; Ho et al., 2020) destroys an image in $T$ small
Gaussian steps and trains a network to predict the noise added, so that sampling can run the
process backwards from pure noise. The \emph{noise schedule} $\beta_1, \dots, \beta_T$ sets how
fast the share of signal left, $\bar\alpha_t = \prod_{s \le t}(1 - \beta_s)$, falls. Ho et al.\
raised $\beta_t$ linearly from $10^{-4}$ to 0.02 over $T = 1000$ steps. Nichol and Dhariwal (2021)
observed that the end of this process is ``too noisy'' to contribute much to sample quality, and
proposed a cosine schedule that destroys information more slowly. Here 32.7\% of
the linear schedule's steps leave less than 1\% of the signal, against 6.5\% for
the cosine schedule, so a sampler that visits only $K$ of the 1000 steps wastes fewer of them on
near-pure noise. Hence the hypothesis: \emph{``A cosine schedule allows a denoising diffusion
probabilistic model to generate high-quality samples using fewer time steps than a linear
schedule.''}

Nichol and Dhariwal found that a linear-schedule model
``does not get much worse'' when up to 20\% of the reverse process is skipped. With the simple
loss used here, their cosine schedule improved log-likelihood on CIFAR-10 (3.37 to 3.26 bits/dim)
but not FID (2.90 to 3.05), and their sampling in far fewer steps came from \emph{learning} the
reverse variances; with fixed variances, as here, quality fell much faster as steps were removed.
Song et al.\ (2021) showed that DDIM, a deterministic sampler for the same trained network, gives
good samples 10 to 50 times faster. Lin et al.\ (2024) showed that common schedules and step
choices leave signal at the last step, which harms sampling in few steps. None of them isolates
what the schedule alone does to few-step sampling, which is the question tested here.

\section{Methodology}
\textbf{Baseline.} Tutorial 9.3's conditional DDPM, unchanged: \texttt{UNet\_cond} (221,569
parameters; its two poolings take $20 \to 10 \to 5$, so $20\times20$ images fit), with $t/T$ and
the class each embedded as 32 numbers added at every level. It predicts the added noise (mean
squared error), trained with Adam ($2\times10^{-4}$) in batches of 128 for 150 epochs on
the 60,000 training images scaled to $[$-$1, 1]$, with $T = 1000$.

\textbf{Cosine schedule.} The only change: $\bar\alpha_t = f(t)/f(0)$ with
$f(t) = \cos^2\!\big(\tfrac{t/T + s}{1 + s}\cdot\tfrac{\pi}{2}\big)$, $s = 0.008$, and
$\beta_t = 1 - \bar\alpha_t/\bar\alpha_{t-1}$ clipped at 0.999 (Nichol and Dhariwal, eq.~17).
Training is otherwise identical: for seeds 0, 1 and 2, both start from the same weights,
batches, steps and noise.

\textbf{Sampling in $K$ steps.} A $K$-step sampler visits $K$ evenly spaced steps rounded to
integers (Nichol and Dhariwal, \S4) and jumps from $t$ to the next visited $t'$ with
$\beta = 1 - \bar\alpha_t/\bar\alpha_{t'}$. Two samplers use it: the tutorial's stochastic step
(DDPM) and the deterministic DDIM step. Both form the clean-image estimate $\hat x_0$ from the
predicted noise and clip it to $[$-$1, 1]$, as Ho et al.'s released sampler does; unclipped at
$K = T$ the step equals the tutorial's to rounding error. The clip matters for the cosine schedule,
whose last $\beta$ is 0.999: the tutorial's form multiplies an error in the first noise estimate by
31.59 at $K = 1000$ and 3,492.67 at $K = 10$ (linear: 2.54).

\section{Experiments}
\textbf{E1} (Figure~1): 12 samples per digit from each seed-0 model at $K = 1000$.
\textbf{E2} (Table~1, Figure~2): each seed-0 model at $K = 10, 20, 50, 100, 250$ with both
samplers and at 1000 with DDPM, 1,000 samples per set, drawn under the labels of a
class-stratified test subset A from the same starting noise for both schedules. \emph{Quality} is
the Fr\'echet distance (FD; Heusel et al., 2017) to a disjoint real subset B with the same class
mix, in the 128 features of a separate classifier (0.9914 test accuracy), as Salvado (2026,
pp.~42, 87) proposes; real A against B gives the floor, 7.65. \emph{Class consistency} is that
classifier's accuracy on the intended labels, and \emph{variety} the within-class feature spread
relative to real digits. \textbf{E3}: seeds 1 and 2 at every $K < 1000$. \textbf{A rule fixed in
advance}: at a given $K$ and sampler, a schedule is better if its FD is lower at seed 0 with a
95\% paired bootstrap interval (1,000 resamples) excluding zero, and lower at seeds 1 and 2 too.
\textbf{E4}: seconds per epoch, and the FD of seed 0's snapshots at epochs 10, 25, 50, 100 and
150 (DDIM, $K = 50$). \textbf{E5}: each model's noise-prediction error on test images at 20
fixed steps. \textbf{E6}, run after the verdict: the seed-0 DDPM runs with Ho et al.'s smaller
posterior variance $\tilde\beta$.

\section{Discussion}
Neither grid in Figure~1 shows mode collapse: every row reads as its digit, with varied slant,
width and stroke (open and closed 4s, looped and plain 2s). The linear model draws even strokes
much like real digits; the cosine model's are heavier, and more of its samples are blotted or
ambiguous (filled 0s, 8s that read as 6s), hence its lower grid score, though on 1,000 samples
the two score alike (0.9650 and 0.9640). The spread agrees: at 1000
steps the cosine samples vary a little more than real digits (1.075 of their
spread) and the linear ones a little less (0.959).

With DDPM sampling, the cosine schedule is better at $K = 10$ and $20$ and neither is clearly better at $K = 50$, $100$ and $250$. With DDIM sampling, neither schedule is clearly better at any $K < 1000$. With DDPM the linear model comes within 10\% of its own 1000-step FD from $K = 100$, the cosine model at no reduced $K$, and the cosine model reaches the linear model's from $K = 100$; with DDIM none of these happens at any reduced $K$. At the full 1000 steps the FDs are 17.12 (linear) and 13.50 (cosine), a difference of $-$3.61 [$-$7.70, $+$0.74], within evaluation noise.

\setcounter{figure}{1}
\begin{figure}[t]
\centering
\includegraphics[width=\columnwidth]{figures/figure_2_steps.pdf}
\caption{FD (top) and class consistency (bottom) against $K$ at seed 0, with seeds 1 and 2 as
small dots; the dotted lines are real test digits.}
\end{figure}

A follow-up at seed 0, run after the verdict, locates the DDPM gain. At $K = 10$ the second-to-last step jumps to step 0 and leaves noise of standard deviation 0.355 (linear) against 0.185 (cosine), which the last step barely removes. With Ho et al.'s smaller posterior variance $\tilde\beta$, the FDs at $K = 10$ go from 186.26 to 42.40 (linear) and from 85.35 to 30.64 (cosine), so most of the cosine schedule's advantage there came from that leftover noise. The cosine schedule thus helps the tutorial's sampler at few steps mostly
through its gentle start near $t = 0$, which shortens that jump, rather than through the gentler
noisy end the hypothesis rests on. DDIM adds no noise, and with it neither schedule is reliably
better: at $K = 10$ seed 0 favours the linear model ($+$57.33) and seed 2 the cosine
($-$4.93). At $K = 10$ the linear model's DDPM samples also lose variety (spread
0.632, cosine 0.944), and the cosine samples are read
as their digit less often at $K = 50$ to 250 with DDPM (0.9430 against
0.9570 at $K = 100$) and at every $K$ with DDIM (0.8180
against 0.9710 at $K = 10$), in line with Figure~1's heavier strokes.

\textbf{Training time and image quality.} An epoch took 7.87~s (linear) and 7.86~s (cosine), the same within timing noise, as expected: the schedule changes which noise level a step uses, not the computation. The seed-0 snapshots reach within 10\% of their final FD by epoch 150 (linear, 19.7~min) and 100 (cosine, 13.1~min). The final training losses differ (0.0283 linear, 0.0492 cosine at seed 0) mostly because the schedules average over different noise levels: 7 of 20 evenly spaced steps leave under 1\% of the signal with the linear schedule and 2 with the cosine, and there the noise is almost free to predict (MSE 0.0007 against 0.0641 elsewhere). At the same noise level with over 1\% of the signal left, the cosine model's error averages 1.00 times the linear model's (0.83 to 1.40), about the same.

\textbf{Strengths and limitations.} The design isolates the schedule (paired weights, batches, noise
and samples), fixes its rule in advance and scales the FD with a real-data floor. Its main
limitation is checkpoint noise: each model is its final checkpoint, without weight averaging, so
its quality moves between epochs (linear FD 44.72,
71.17 and 28.48 at epochs 50, 100 and 150) and between seeds
(cosine, DDPM, $K = 250$: 46.43 at seed 1, 15.58 at seed 0).
That spread, not the sample count, leaves most step counts unclear and makes the epoch figures above
a single noisy reading. The intervals cover seed 0's sampling only; the FD uses a small classifier's
features at 1,000 samples, so it is not comparable with FID; and every model is a
$T = 1000$ U-Net on $20\times20$ digits.

\textbf{Recommendations.} Average the weights during training, as Ho et al.\ do (an exponential
moving average, decay 0.9999), and add seeds. Sample in few steps with the posterior variance $\tilde\beta$ or with learned variances,
which Nichol and Dhariwal found keep quality with an order of magnitude fewer steps. To test a
smaller $T$, rescale the linear $\beta$ range with $T$, as their released code does; the tutorial's
ends at $\bar\alpha_T = 0.364$ over 100 steps.

\section{Conclusion}
The evidence supports the hypothesis in part: by a rule fixed before the results, the cosine schedule gave better samples at $K = 10$ and $20$ with DDPM, and the linear schedule at no reduced step count. Most of that gain comes from the noise the tutorial's sampler leaves after its last coarse step,
which the posterior variance removes; with DDIM neither schedule is reliably better. A stronger
claim needs averaged weights and more seeds.

\section*{References}
\footnotesize
\setlength{\parindent}{0pt}
\newcommand{\bib}[1]{\par\hangindent=1em #1}
\bib{Heusel, M. et al. (2017). GANs trained by a two time-scale update rule converge to a local
Nash equilibrium. \emph{NeurIPS}.}
\bib{Ho, J., Jain, A. and Abbeel, P. (2020). Denoising diffusion probabilistic models.
\emph{NeurIPS}.}
\bib{Lin, S., Liu, B., Li, J. and Yang, X. (2024). Common diffusion noise schedules and sample
steps are flawed. \emph{WACV}.}
\bib{Nichol, A. and Dhariwal, P. (2021). Improved denoising diffusion probabilistic models.
\emph{ICML}.}
\bib{Salvado, O. (2026). Week 9: transforming distribution (GAN, diffusion, flow). IFN680
lecture, QUT.}
\bib{Song, J., Meng, C. and Ermon, S. (2021). Denoising diffusion implicit models. \emph{ICLR}.}

\end{document}
